% !TITLE 观刈麦
% !AUTHOR 白居易
% !DYNASTY 唐
% !GENRE 五言古诗
% !ORDER 45

\begin{poem}
田家少闲月，五月人倍忙。
夜来南风起，小麦覆陇黄\textsuperscript{1}。
妇姑荷箪食\textsuperscript{2}，童稚携壶浆\textsuperscript{3}，
相随饷田\textsuperscript{4}去，丁壮在南冈。
足蒸暑土气，背灼\textsuperscript{5}炎天光，
力尽不知热，但惜夏日长\textsuperscript{6}。
复有贫妇人，抱子在其旁，
右手秉遗穗\textsuperscript{7}，左臂悬敝筐\textsuperscript{8}。
听其相顾言，闻者为悲伤。
家田输税尽\textsuperscript{9}，拾此充饥肠。
今我何功德？曾不事农桑。
吏禄三百石\textsuperscript{10}，岁晏有余粮\textsuperscript{11}，
念此私自愧，尽日不能忘。
\end{poem}

\begin{annotations}
\notetext{1}{覆陇黄：麦子覆盖田垄，一片金黄。覆，覆盖；陇，田垄。}
\notetext{2}{荷箪食：挑着竹器装的饭食。荷（hè），挑；箪（dān），圆形竹器。}
\notetext{3}{携壶浆：提着壶里的汤水。浆，汤水或米汤。}
\notetext{4}{饷田：到田里去送饭。饷（xiǎng），给人送食物。}
\notetext{5}{灼：烤灼，烈日照在背上如火烤一般。}
\notetext{6}{但惜夏日长：只担心夏天太短，天黑了就干不成活。不是不热，是顾不上热——农人的艰辛在这种反常逻辑中呈现。}
\notetext{7}{秉遗穗：拿着别人收割后遗漏的麦穗。秉（bǐng），持、拿；遗穗，掉落的麦穗。}
\notetext{8}{悬敝筐：臂上挂着破旧的竹筐。敝（bì），破旧。}
\notetext{9}{家田输税尽：家里的田地因为交税都卖光了。输税，缴纳赋税。}
\notetext{10}{吏禄三百石：县尉的年俸禄是三百石粮食。石（dàn），容量单位。}
\notetext{11}{岁晏有余粮：年底还有剩余的粮食。晏（yàn），晚、末。白居易对比自己不劳而获与农人的劳而不获，发自真心地愧疚。}
\end{annotations}

\begin{appreciation}
刈是割的意思。观刈麦，就是站在田边看人割麦。这是白居易早年的名篇之一，也是中国文学史上把”悯农”写得最不自欺的诗之一——写农人辛苦的诗很多，这一首特别在：作者把自己也写了进去。

田家少闲月，五月人倍忙——庄户人家一年到头没有闲月，五月倍忙。夜来南风起，小麦覆陇黄——夜里南风一起，次日田垄全黄。麦熟就要抢收，一场雨就能让一年的辛苦泡汤——金黄不是风景，是催命符。

妇姑荷箪食，童稚携壶浆，相随饷田去——婆婆媳妇挑饭，小娃娃提水，送到田里去。饷田是送饭到田里，箪是竹筐，浆是米汤。全家老小，没有闲人。

足蒸暑土气，背灼炎天光——脚下蒸着暑气，背上烤着日头。力尽不知热，但惜夏日长——力气用尽了，反倒不觉得热，只恨天黑太早。这两句是诗眼：不是不热，是热到顾不上；不是不怕累，是怕活干不完。白居易不解释，让反常说话。

复有贫妇人，抱子在其旁，右手秉遗穗，左臂悬敝筐——一个抱着孩子的穷妇人，在田边捡别人落下的麦穗。她家的地早卖了交税。家田输税尽，拾此充饥肠——捡麦穗不是为了省，是为了活命。

诗本可在这里收尾。可白居易没有停，他忽然问自己：今我何功德？曾不事农桑。吏禄三百石，岁晏有余粮——我不种地不养蚕，凭什么一年三百石俸禄、年底有余粮？念此私自愧，尽日不能忘——这一想，他惭愧了一整天。

这一问，把诗从”看别人”变成了”照自己”——后来新乐府主张”为事而作”，写诗的人，先把自己摆进去照一照。这本书后面还有他的《卖炭翁》——他一辈子都在让城里人看见粮食和炭火的来处。《氓》里那个女子说”自我徂尔，三岁食贫”——穷人的日子，诗人们写了几千年也没写完。

你碗里的每一粒米饭，都是从这样的五月里来的。别的可以不记，这一条记住就行：好好吃饭，别剩饭。
\end{appreciation}
